\chapter{Теоретические основы блокчейн-технологий}\label{ch:ch1}

\section{Исторические предпосылки развития блокчейна}\label{sec:ch1/sec1}

Блокчейн как технология является результатом накопления знаний и достижений в
различных научных областях. Его создание стало возможным благодаря эволюции
криптографии, теории распределенных систем, теории игр и экономических моделей
стимулирования участников. Эти дисциплины сформировали теоретические и
практические основы, позволившие решать ключевые задачи распределенных систем:
безопасность, консенсус, масштабируемость и мотивацию участников.

История блокчейна начинается задолго до появления Биткойна в 2008 году. Корни
уходят в середину XX века, когда были заложены основы криптографических методов,
концепций консенсуса и протоколов распределенных систем. Ниже рассмотрены
основные направления, которые сформировали технологическую базу для блокчейна.

\subsection{Криптография}\label{sec:ch1/sec1/crypto}

Криптография обеспечивает безопасность, неизменность данных и доверие между
участниками сети. Ее современное развитие началось в XX веке, когда защита
информации стала критически важной задачей. Разработка симметричных алгоритмов,
таких как DES, дала основу для защиты данных, но потребовала безопасного обмена
секретными ключами. Прорывом стала концепция открытых и закрытых ключей,
предложенная Диффи и Хеллманом, а затем практическая реализация RSA, которая
заложила фундамент цифровых подписей~\cite{DiffieHellman1976,Johnson1994}.

Важным этапом стало развитие криптографических хэш-функций. Их свойство
необратимости и чувствительности к изменениям делает возможным контроль
целостности данных. Концепция хэш-деревьев Меркла позволила эффективно проверять
большие объемы данных без их полного раскрытия, что легло в основу структуры
блокчейна~\cite{Merkle1980}.

Отдельное значение имеет концепция Proof-of-Work, впервые предложенная как защита
от спама и впоследствии адаптированная для консенсуса в сети Биткойн. Наличие
вычислительной сложности делает атаки на сеть экономически невыгодными, а данные
устойчивыми к подделке~\cite{Back2002,Nakamoto2008}.

\subsection{Распределенные системы}\label{sec:ch1/sec1/dist}

Распределенные системы исследуют надежность и согласованность данных в условиях
отсутствия единого центра. Классические работы по проблеме византийских
генералов и протоколы консенсуса (Paxos, Raft) сформировали основу для
механизмов согласования состояния между независимыми узлами.
~\cite{Lamport1982,Lamport2001,Ongaro2014}

Блокчейн использует эти идеи, но дополняет их экономическими стимулами и
криптографическими механизмами. Это позволяет обеспечить согласованность данных
в открытых сетях, где участники не доверяют друг другу и могут действовать
злонамеренно. Развитие распределенных систем также повлияло на дизайн сетей с
разделением нагрузки (шардинг), что стало важным направлением масштабирования.

\subsection{Теория игр}\label{sec:ch1/sec1/game}

Теория игр внесла вклад в формирование механизмов экономической мотивации и
стимулирования честного поведения. Блокчейн-системы используют модели, в которых
участники принимают решения, исходя из индивидуальной выгоды. Консенсусные
алгоритмы проектируются так, чтобы честное поведение было рациональной стратегией
для большинства участников~\cite{VonNeumann1944,Eyal2014}.

В частности, модели большинства и устойчивости к коалициям участников позволяют
анализировать риски централизации и атак, таких как 51\%-атака. Эти подходы
используются при проектировании и анализе протоколов PoW и PoS.

\subsection{Экономика}\label{sec:ch1/sec1/econ}

Экономические модели лежат в основе токеномики и механизмов стимулирования. Для
блокчейн-сетей критично обеспечить баланс между вознаграждением валидаторов,
стоимостью участия и устойчивостью сети. Эмитируемые токены выполняют роль
стимулов за поддержание инфраструктуры и валидность транзакций.

Развитие децентрализованных финансов и механизмов управления активами
подтвердило, что экономические инструменты являются неотъемлемой частью
архитектуры блокчейна. Это особенно важно для сетей, ориентированных на
автоматизацию инфраструктуры валидаторов, где экономические стимулы определяют
качество и надежность предоставляемых ресурсов~\cite{MakerDAO2017,Uniswap2018}.

\section{Основные компоненты и принципы работы блокчейна}\label{sec:ch1/sec2}

Блокчейн представляет собой распределенный реестр, в котором данные хранятся в
цепочке последовательных блоков, связанных криптографическими хэшами. Каждая
запись вносится по правилам консенсуса, а изменение уже включенных данных
становится практически невозможным без пересчета всей цепочки.

Ключевые компоненты блокчейн-системы:
\begin{itemize}
  \item модель данных и структура блоков, обеспечивающие неизменность записей;
  \item механизмы консенсуса, определяющие правила подтверждения транзакций;
  \item роль валидаторов и распределение ответственности между участниками;
  \item сетевой уровень, обеспечивающий распространение данных и устойчивость к
  отказам.
\end{itemize}

Эти принципы формируют базу для дальнейшего рассмотрения архитектуры валидаторской
инфраструктуры и требований к автоматизации в последующих главах.
